% !TEX root =../main.tex
\section{Marco Teórico}

\subsection{Modelo OSI}
  El modelo OSI(Open Systems Interconnection) es un modelo de referencia\cite{ISO_OSI} 
  creado por International Organization for Standardization(ISO) para la 
  comunicación de diversos sistemas computacionales\cite{cloudflare_OSI}, 
  se encuentra dividido en 7 capas que describen funciones especificas 
  y proveen servicios a las capas superiores.

  \paragraph{1. Capa Física}
  La capa física tiene como función especificar y describir las caracteristicas 
  fisica de los medios y las conexiones usadas para la tranmisión de información
  como señales electricas o electromagneticas(opticos, inalambricos).
  Para cada medio se describen sus caracteristicas mecanicas y electricas, el como codifica la información,
  su ancho de banda y distancias maximas.\cite{netacad_physical_layer}
  En esta capa se pueden ubicar los estandares Ethernet, WLAN(WirelessLAN), medios como cables UTP, fibra óptica y
  conectores como el RJ45
  \paragraph{2. Capa de Enlace de Datos}
  Tiene como objetivo proporcionar transferencia de datos confiable entre dos dispositivos 
  conectados físicamente. Se encarga de corregir errores ocurridos en la capa física y provee
  acceso a esta para las capas superiores. Además incluye un direccionamiento físico para el envio
  de los datos entre dispositivos.\cite{osi_windows}
  Puede dividirse en subcapas MAC(Media Access Control) y LLC(Logical Link Control)
  que se encargan de : el control del medio físico, identificandolo con una dirección y corrigiendo errores;
  y la transferencia de información sin errores, control de flujo para los datos y la retransmisión.
  En esta capa los datos(PDU) se denominan "frames" que contienen la información enviada desde la capa de redes
  más el direccionamiento MAC.

  \paragraph{3. Capa de Red}
  Es la capa encargada de la comunicación entre diferentes redes, se encarga encarga de el ruteo
  para enviar un mensaje desde el host productor hasta el host destino, esto mediante la dirección 
  lógica asignada a cada equipo. Utiliza datos(PDU) denominados "packets" que contiene información
  de capa 4 más direccionamiento IP(Internet Protocol).\cite{cloudflare_OSI}
  \paragraph{4. Capa de Transporte}
  La capa de transporte se encarga de segmentar la información y secuenciarla para el para luego enviarla
  en capas inferiores, además de recomponerla cuando la reciba, implementa control de flujo para regular
  la velocidad de transmisión y corrección de errores para asegurar la comunicación de la información.\cite{osi_aws}
  Los protocolos TCP(Transmission Control Protocol) y UDP(User Datagram Protocol) son los principales
  actores en la capa de transporte.
  \paragraph{5. Capa de Sesión}
  La capa de sesión se encarga de establecer, abrir y cerrar comunicaciones entre maquinas,
  maneja la sincronización para contenter interrupciones en la conexión y control de tráfico.\cite{osi_aws}
  \paragraph{6. Capa de Presentación}
  Capa encargada de la traducción, encriptación y compresión de datos, para ser enviados capa arriba o capa abajo.
  Los datos recibidos desde la capa 5 son convertidos a formatos digeribles por la capa de aplicación como lo son HTML, CSV; 
  y los datos recibidos desde la capa 7 son comprimidos y encriptados antes de ser enviados a la capa 5.
  \paragraph{7. Capa de Aplicación}
  Es la capa que tiene interacción directa con el usuario incluye aplicaciones de software como los buscadores,
  transferencia de archivos, conexión remota, llamadas, streaming.
  Algunos protocolos utilizados en esta capa son 
  HTTPS(Hypertext Transfer Protocol Secure), SSH(Secure Shell) o FTP(File Transfer Prolocol)\cite{osi_Haryana_India}
\subsection{Dispositivos de networking}
  \paragraph{Router}
  Un router es un dispositivo encargado de enviar mensajes de una red a otra, esto lo hace con funciones de capa 3,
  leyendo la dirección IP contenida en el "packet" y enviando la información hacia el siguiente dispositivo que permita
  al mensaje llegar a su destino, mediante funciones de capa 2. Se comunican mediante Ethernet, fibra optica o WLAN y necesitan ser configurados antes de ser
  utilizados, tienen su propia dirección IP la cual usan para ser identificados por otros dispositvos.\cite{cloudflare_router}
  \paragraph{Switch}
  Un switch es un dispositivo encargado de enviar mensajes al interior de una red, utiliza funciones de capa 2 para enviar mensajes entre
  dispositivos, mediante el uso de el direccionamiento MAC, el switch conecta directamente 2 dispositivos de red manteniendo el ancho de banda.\cite{cloudflare_switch}
  Un switch se comunica mediante Ethernet o fibra optica, no requiere de configuración para hacer su trabajo 
  sin embargo si requiere hacer más tareas, como monitoreo, debe configurarse según se requiera.
  \paragraph{Hub}
  Un hub es un dispositivo que replica el trafico que recibe, reenviando los mensajes a todos los dispositivos conectados a el,
  utiliza unicamente funciones de capa 1 por lo que se puede considerar como un cable más.
  \paragraph{Cliente o Computador}
  \paragraph{Firewall}
